%% Section 1: Introduction
\section{Introduction}
\label{sec:intro}
\IEEEPARstart{G}{rid-connected} photovoltaic (PV) generation is expanding rapidly, and the single-phase inverter is the component that largely determines conversion efficiency, cost and power quality at residential and commercial scale \cite{Kjaer2005}. Line-frequency or high-frequency transformers provide galvanic isolation between the PV array and the grid, but they add weight, cost and loss. Removing the transformer improves efficiency and power density \cite{Gonzalez2007,Araujo2010}, and transformerless inverters are therefore widely used in the low-power range.

The price of removing isolation is a conductive path between the PV array and the grid ground. The PV modules have a parasitic capacitance $\Cpv$ to the grounded frame and mounting structure, and this capacitance is in series with the inverter output filter and the grid impedance. Any time-varying CMV generated by the inverter drives a leakage current through $\Cpv$ \cite{Gubia2007,Lopez2007}. Such a current causes electromagnetic interference, additional loss, grid-current distortion and, above all, a shock hazard. Safety standards require disconnection when the continuous residual current exceeds 300~mA, and when it changes suddenly by specified steps within specified times \cite{VDE0126,IEC62109}.

The classical remedy is topological: H5, H6 and HERIC-type inverters decouple the PV array from the grid during the freewheeling intervals so that the CMV remains (ideally) constant \cite{Victor2008,Schmidt2003,Zhang2014,Kerekes2011}. In practice, however, the achieved leakage current depends on non-idealities and on the operating point. The switching frequency is of particular interest: it sets the filter size and the switching loss, but it also determines where the harmonics of the CMV fall relative to the CM resonance of the $L$--$\Cpv$ loop. This interaction is usually discussed qualitatively, and switching frequency is commonly chosen on the basis of efficiency and filter size alone.

\subsection{Contributions}
This paper makes the following contributions.
\begin{enumerate}
  \item A frequency-resolved simulation study (1--40~kHz) of the leakage current, efficiency and low-order THD of the H4 inverter with unipolar and bipolar sinusoidal PWM (SPWM) and of the HERIC inverter with unipolar and bipolar SPWM, using an L filter and a 150~nF parasitic capacitance.
  \item A CM-resonance interpretation of the results: in all four cases the leakage peaks occur where $\fsw$ coincides with a low-order sub-multiple of $\fres$ (or, in one case, with $2\fres$), where $\fres=1/(2\pi\sqrt{\Leq\Cpv})$ evaluated from the circuit parameters, and the leakage floor of the HERIC bipolar case is consistent with the line-frequency term $v_o/2$ of the effective CM source acting on $\Cpv$.
  \item A practical switching-frequency selection rule, together with a comparison of the operating windows that satisfy the 300~mA limit, and an explicit account of the limitations of the study (simulation only, non-uniform parameters between H4 and HERIC cases, band-limited THD).
\end{enumerate}

\subsection{Organization}
Section~\ref{sec:related} reviews related work and positions the study. Section~\ref{sec:cm} derives the CM model and the resonance condition. Section~\ref{sec:topologies} describes the topologies and modulation schemes. Section~\ref{sec:method} documents the simulation set-up and the metrics. Section~\ref{sec:results} presents the results, Section~\ref{sec:discussion} discusses the results, Section~\ref{sec:design} gives the sensitivity analysis and design guidelines, Section~\ref{sec:limits} states the limitations, and Section~\ref{sec:conclusion} concludes.
